% Part: lambda-calculus
% Chapter: syntax
% Section: abbreviated-syntax

\documentclass[../../../include/open-logic-section]{subfiles}

\begin{document}

\olfileid{lam}{syn}{abb}
\olsection{సంక్షిప్త వాక్యనిర్మాణం}

\olref[trm]{defn:term}లో నిర్వచించిన విధంగా పదాలను రాయడం కొన్నిసార్లు
క్లిష్టమవుతుంది; అందుకే మరింత సంక్షిప్త వాక్యనిర్మాణాన్ని ప్రవేశపెట్టడం
ఉపయోగకరం. సంక్షిప్త సంకేతనంలోని పదాలకూ ఏకార్థ పఠనం ఉండేలా జాగ్రత్త
పడాలి. విస్తృతంగా వాడే ఒక రూపాన్ని \emph{సంక్షిప్త పదాలు} అంటారు.
దాని నియమాలు ఇవి:

\begin{enumerate}
\item కుండలీకరణలను వదిలినప్పుడు ప్రయోగం ఎడమ నుంచి కుడికి జరుగుతుంది.
  ఉదాహరణకు, $M$, $N$, $P$,~$Q$ పదాలైతే, $MNPQ$ అనేది
  $(((MN)P)Q)$కు సంక్షిప్త రూపం.
\item కుండలీకరణలను వదిలినప్పుడు లాంబ్డా అమూర్తీకరణకు సాధ్యమైనంత
  విస్తృత పరిధి ఇస్తాం. ఉదాహరణకు, $\lambd[x][MNP]$ను
  $(\lambd[x][MNP])$గా చదువుతాం.
  \sourcecorrection{OLTELAMSYN-004}{ఆంగ్ల మూలం ఇక్కడ ``From example'' అని రాసింది; సందర్భానికి తగిన ``For example'' అర్థాన్ని అనువదించాం.}
\item ఒకే లాంబ్డాతో అనేక చరాలను అమూర్తీకరించవచ్చు. ఉదాహరణకు,
  $\lambd[xyz][M]$ అనేది $\lambd[x][\lambd[y][\lambd[z][M]]]$కు
  సంక్షిప్త రూపం.
\end{enumerate}

ఉదాహరణకు,
\[
\lambd[xy][xxyx \lambd[z][xz]]
\]
అనేది
\[
(\lambd[x][(\lambd[y][((((xx)y)x)(\lambd[z][(xz)]))])]).
\]
కు సంక్షిప్త రూపం.

\begin{prob}
$\lambd[g][(\lambd[x][g (x x)])
  \lambd[x][g (x x)]]$ అనే సంక్షిప్త పదాన్ని పూర్తి రూపంలో రాయండి.
\end{prob}

\end{document}
